% 本文件由 LaTeX 排版 Agent 根据有效 translation.json 自动生成。
% author=Malchion; work_id=0617; title=书信

\chapter{书信}

% structure: p0001 source=h1 target=omitted reason=page-title-already-rendered-by-parent

\Emph{以安提约基雅会议之名，反对萨摩萨塔的保禄。}\par

致\Person{狄约尼削}与\Person{玛克西穆}，并致全世界所有与我们一同在牧职中的同道，包括主教、长老与执事，并致普天之下整个天主教会：\Person{赫肋诺}、\Person{依默纳约}、\Person{德敖斐罗}、\Person{德敖德诺}、\Person{玛克西穆}、\Person{普洛科鲁}、\Person{尼各玛}、\Person{厄里亚诺}、\Person{保禄}、\Person{波拉诺}、\Person{普洛托革诺}、\Person{耶辣克}、\Person{欧提基}、\Person{德敖多禄}、\Person{马尔基翁}、\Person{路济亚}，以及所有与我们同在、居住在邻近城市与民族中的其他主教、长老与执事，连同天主的众教会，向我们在主内亲爱的弟兄姊妹致候。\par

1. 在简短的开场白之后，他们这样继续说道：——我们写信给许多主教，甚至那些住在远处的，并劝他们帮助我们摆脱这致命的教义；在这些主教中，我们特别写信给亚历山大里亚的主教\Person{狄约尼削}和卡帕多细亚的\Person{菲尔米里安}，那些有福名声的人。其中，一位写信给安提约基雅，甚至不屑于以他自己的名义致信这位错误的领袖；他不是以他自己的名义写信给他，而是写给整个地区。我们也附上了那封信的副本。而\Person{菲尔米里安}，他曾两次亲自前来，谴责了教义上的革新，我们这些在场的人知道并作证，许多其他人也和我们一样知道。但是当保禄承诺放弃这些意见时，他相信了他；并希望在没有对圣言造成任何羞辱的情况下，事情能得到妥善解决，他推迟了决定；然而，在这件事上，他被那个否认他的天主和主、并背叛他先前所持信仰的人欺骗了。现在\Person{菲尔米里安}打算前往安提约基雅；他到达了\Place{塔尔索}，仿佛已经对那人的无信邪恶有所体验。但当我们刚刚聚集并呼唤他、等待他的到来时，他的结局临到了他身上。\par

2. 在谈论其他事情之后，他们以下列方式告诉我们他的生活方式：——但是，当他还在教会之外时，当他背离信仰并转向虚假和非法的教义时，没有必要评判他的行为。我们也不需要说任何这类的事情，例如：他以前贫穷可怜，没有从父辈继承任何财产，也没有通过技艺或任何行业获得任何财富，但现在他通过不义和亵渎的行为，以及通过那些他掠夺和勒索弟兄们的手段，拥有了过度的财富，他在诉讼中不公平地对待受伤害者，并承诺以代价帮助他们，却一直欺骗他们并使他们损失，利用那些处于困境中的人们愿意付出以从困扰中解脱的意愿，从而以为敬虔便是牟利\BibleRef{弟前 6:5}。我也不需要说任何关于他的骄傲和傲慢，他以之掌握世俗的尊荣，并希望被称为总督而非主教，以及他在市场里高视阔步，公开边走边读并朗诵信件，并被众多人群前呼后拥；以致他以傲慢的姿态和内心的狂妄为信仰招致恶意和仇恨。我也不说他在教会集会中施行的骗术，以追求声望、大张旗鼓，并用这种手腕使头脑较单纯的人惊异；也不说他为自己设立了一个高高的讲台和宝座，如此不像基督的门徒；也不说他有一个\OriginalTerm{秘密室}{secretum}并以此名称呼它，仿效世界统治者的方式；也不说他用手拍打大腿，用脚踢打讲台；也不说他谴责和侮辱那些不鼓掌、不摇手帕（如同在剧院中所为）、不像他的男男女女支持者那样大喊大叫、跳跃蹦跶——这些支持者对他是如此无秩序的听众，却选择像在天主家中一样恭敬谦卑地聆听——的人；也不说他在会众中对那些已经去世的圣言的讲解者进行不当和猛烈的攻击，并高抬自己，不像一位主教，倒像一位诡辩家和魔术师；也不说他停止那些为赞美我们的主耶稣基督而唱的圣咏，声称是新近的人的新近作品，并准备妇女们在教会中间，在逾越节大日，唱赞美他自己的圣咏——这样的唱诗班让人一听就发抖。此外，他还影响那些在邻近地区和城市中奉承他的主教和长老们，在他们的讲道中向百姓推进类似的意见。\par

3. 我们可以说——为了稍作预告，我们打算在下面写些什么——他不愿意承认天主圣子从天降下。这个说法并非仅凭断言；因为我们寄给你们的那些备忘录充分证明了这一点，尤其是那段他说耶稣基督是来自下面的经文。那些在民众中歌颂他、赞美他的人，宣称他们那不敬神的老师像天使一样从天降下。这个傲慢的人并不制止这些言论，甚至当这些言论发出时他也在场。此外，还有这些妇女——安提约基雅人称之为\Quote{义姊妹}——由他以及与他在一起的司铎和执事们供养，尽管他知道并已证实她们在此事及其他事上犯有不可救药的罪过，他却纵容隐瞒，目的是使那些男人顺服于自己，并阻止他们因自身的地位而害怕，从而不敢胆敢就他那不敬神的言行指控他。除此之外，他使他的追随者富有，因此那些贪恋这些东西的人爱慕他、钦佩他。但我们为什么要写这些事呢？因为，亲爱的，我们知道主教和全体圣职人员应该在一切善行上成为民众的榜样。我们并非不知道许多人因将这类妇女引入家中而堕落，另一些人则因此受嫌疑。所以，即使有人承认他在这件事上没有做出什么可耻之事，他至少也应该避免由这种行径所产生的嫌疑，以免有人被冒犯，或找到模仿他的理由。因为，这样看来，又岂能有人指责别人，或警告他要提防与女子过于亲近，以免失足呢？正如经上所记：\BibleQuote{如果一个人虽然辞退了一个，却仍留有两个在他身边，且她们正值青春年少、容貌美丽；而且他外出时又带着她们同行；同时他又过着奢侈放纵的生活？}\par

4. 因这些事，众人都私下叹息哀怨；然而他们对他的暴虐与威权如此惧怕，以致不敢冒险指控他。对于这些事，正如我们已经说过，如果是属于我们中间、持有公教情怀的人，或许会予以计较；但对于一个已经背叛了（信仰的）奥迹，并与可憎的阿尔特玛斯异端一道夸口的人——我们为何要迟疑不揭发他的根源呢？——我们认为无需就这些事向他追究。\par

5. \Emph{接着，在书信的结尾，他们添加了以下话语}：——因此，我们被迫开除这个反对天主自身且拒不服从的人，并为他选立另一位主教来牧养天主教会，这是依靠天主的眷顾——即德梅特里亚努斯之子，一位蒙福记忆者，他曾杰出地治理过同一教会，名叫多米努斯，一个具备主教所应有的一切高贵品质的人。并且我们已经将此事实通知你们，以便你们可以写信给他，并从他那里接受共融书信。而那另一位若愿意，可以写信给阿尔特玛斯；凡与阿尔特玛斯思想相同的人，若愿意，也可以与他共融。\par

\SourceRecord{Translated by S.D.F. Salmond.；Ante-Nicene Fathers；Vol. 6.；Edited by Alexander Roberts, James Donaldson, and A. Cleveland Coxe.；Buffalo, NY: Christian Literature Publishing Co.,；1886.；<http://www.newadvent.org/fathers/0617.htm>.}
